\documentclass[11pt,a4paper]{article}
\usepackage{mathtools}
\usepackage[utf8]{inputenc}
\usepackage[T1]{fontenc}
\usepackage[british]{babel}

% ── Mathematics ───────────────────────────────────────────────────────────────
\usepackage{amsmath,amssymb,amsthm}
\usepackage{mathtools}

% ── Layout ────────────────────────────────────────────────────────────────────
\usepackage[margin=2.5cm]{geometry}
\usepackage{setspace}
\onehalfspacing
\usepackage{booktabs}
\usepackage{array}
\usepackage{enumitem}
\usepackage[hidelinks]{hyperref}

% ── Bibliography ──────────────────────────────────────────────────────────────
\usepackage[numbers,sort&compress]{natbib}
\bibliographystyle{plainnat}

% ── Theorem environments ──────────────────────────────────────────────────────
\theoremstyle{plain}
\newtheorem{theorem}{Theorem}[section]
\newtheorem{proposition}[theorem]{Proposition}
\newtheorem{lemma}[theorem]{Lemma}
\newtheorem{corollary}[theorem]{Corollary}
\theoremstyle{definition}
\newtheorem{definition}[theorem]{Definition}
\theoremstyle{remark}
\newtheorem{remark}[theorem]{Remark}
\newtheorem{openprob}{Open Problem}

% ── PDL macros ────────────────────────────────────────────────────────────────
\newcommand{\PDL}{\textsc{PDL}}
\newcommand{\Kfour}{K_4}
\newcommand{\rval}{r_{\mathrm{val}}}
\newcommand{\Rsea}{R_{\mathrm{sea}}}
\newcommand{\Rtot}{R_{\mathrm{tot}}}
\newcommand{\Rsurf}{R_{\mathrm{surf}}}
\newcommand{\Gval}{G_{\mathrm{val}}}
\newcommand{\Gsea}{G_{\mathrm{sea}}}
\newcommand{\vp}{\varphi}
\newcommand{\kap}{\kappa}
\newcommand{\sig}{\sigma}
\newcommand{\GPDL}{G_{\mathrm{PDL}}}
\newcommand{\Geff}{G_{\mathrm{eff}}}
\newcommand{\hb}{\hbar}
\newcommand{\SBH}{S_{\mathrm{BH}}}
\newcommand{\SPDL}{S_{\mathrm{PDL}}}
\newcommand{\Omsurf}{\Omega_{\mathrm{surf}}}
\newcommand{\lP}{l_P}

% ── Title ─────────────────────────────────────────────────────────────────────
\title{\textbf{Surface Locality of Relational Entropy\\
and the PDL Area Law}\\[4pt]
\large Document D37 --- PDL Programme}
\author{C\'edric Laubscher\\
\small Independent Researcher, Switzerland}
\date{March 2026}

% ══════════════════════════════════════════════════════════════════════════════
\begin{document}
\maketitle

% ── Abstract ──────────────────────────────────────────────────────────────────
\begin{abstract}
Within the Projective Dynamic Logo (\PDL) framework, every composite structure is
organised as a relational closure whose interior (the valence sector, $\rval = 930$
cycles) communicates with the exterior exclusively through a finite active surface
($\Rsurf = 310\vp$ cycles, $\vp$ the golden ratio). Paper~D29 proved that the valence
sector occupies a \emph{unique} configuration satisfying the $(A)\wedge(B)$ stability
criterion; its accessible microstate count is therefore exactly one.
We formalise this observation as \emph{Proposition~BH-1}: the information-theoretic
entropy of a \PDL\ closure is carried entirely by the surface sector, not by the
total relational budget. This constitutes an unconditional structural derivation of the
\emph{area law} $S \propto \Rsurf$, which maps --- via the Gate~3 theorem
$\Geff(N)=\sig(N)\cdot\GPDL$ (D31/D36) --- to the Bekenstein--Hawking scaling
$\SBH \propto A$ at macroscopic scale. The coefficient $1/4$ in the Bekenstein--Hawking
formula is identified as a quantitative open problem (OP1) conditional on the algebraic
proof of~$\kap = \Rsurf/\Rtot$ (OP1 of D36). Epistemic status: BH-1 is an unconditional
theorem within the \PDL\ corpus; its macroscopic interpretation involves the Gate~3
hypothesis H1.
\end{abstract}

\newpage

\tableofcontents

\newpage

% ══════════════════════════════════════════════════════════════════════════════
\section{Introduction and motivation}
\label{sec:intro}

The Bekenstein--Hawking entropy formula
\begin{equation}
  \SBH = \frac{k_B c^3 A}{4\,G\,\hb}
  \label{eq:BH}
\end{equation}
is one of the most constrained results in theoretical physics: it is supported by
Euclidean path-integral arguments, string theory microstate counting, and the
membrane paradigm, yet its derivation always presupposes curved spacetime and quantum
fields on that spacetime. The deep reason why entropy is proportional to \emph{area}
rather than \emph{volume} remains an axiom-level postulate in standard approaches,
adopted under the name of the holographic principle.

The \PDL\ framework derives physics from four axioms on finite signed graphs, without
presupposing spacetime, fields, or thermodynamic postulates. Its central object, the
proton, is characterised by the integer quintuplet $(24,28,930,10087,11017)$ and
exhibits a precise internal architecture: a \emph{valence sector} of $\rval=930$
stationary relations and a \emph{sea sector} of $\Rsea=10087$ dynamical relations,
coupled to an external $\Kfour$ electron prototype exclusively through an \emph{active
surface} $\Rsurf = 310\vp \in \mathbb{Q}(\sqrt{5})$.

The present paper asks a sharp question: \emph{how many externally accessible
microstates does a \PDL\ closure admit?} The answer, derived from existing \PDL\
theorems, is that the valence sector contributes exactly one accessible microstate
(its configuration is uniquely fixed by the $(A)\wedge(B)$ criterion, proved in D29),
so that the entropy is entirely carried by the surface sector. This is BH-1, stated
and proved in Section~\ref{sec:bh1}. Section~\ref{sec:macro} connects BH-1 to
Bekenstein--Hawking via the Gate~3 theorem. Section~\ref{sec:open} identifies the
open quantitative problem (the $1/4$ coefficient). Section~\ref{sec:epistemic} gives
the epistemic status table.

\paragraph{Relation to prior work.}
The structural analogy between the proton and a black hole was introduced qualitatively
in the main \PDL\ paper~\cite{PDL_Main_2026} and developed as the ``inverted mini
black hole'' picture. The present paper is the first to give a precise, theorem-level
formulation of the area law within \PDL, grounded in the results of D29--D36.

% ══════════════════════════════════════════════════════════════════════════════
\section{Prerequisites}
\label{sec:prereq}

We recall the \PDL\ results used below. All proofs are in the cited papers.

\subsection{The proton quintuplet and its partition}

The proton is characterised by
$(n_u, n_d, \rval, \Rsea, \Rtot) = (24, 28, 930, 10087, 11017)$,
where $n_u, n_d$ are the vertex counts of the two valence cores,
$\rval = 2\binom{n_u}{2} + \binom{n_d}{2} = 930$ their stationary relation count,
$\Rsea = 10087$ the dynamical sea relations, and $\Rtot = \rval + \Rsea$
\cite{PDL_Main_2026,Combinatorial_Proton_PDL_2026}.

\subsection{The \texorpdfstring{$(A)\wedge(B)$}{(A)∧(B)} criterion and valence uniqueness}

\begin{definition}[$(A)\wedge(B)$ stability, D29~\cite{PDL_D29_2026}]
A mixed triangle $(e_i,e_j,p_k)$ is \emph{stable} if and only if the cross-product
$P_c = s_\times(e_i,p_k,c)\cdot s_\times(e_j,p_k,c)$ satisfies:
\begin{align}
  (A) &:\quad P_1 = s_{K_4}^{(1)}(e_i,e_j), \\
  (B) &:\quad P_2 = -P_1.
\end{align}
Equivalence $(A)\wedge(B) \iff \text{stable}$ was verified exhaustively over $768$
cases with zero counter-example.
\end{definition}

\begin{theorem}[Valence uniqueness, D29 Lemma~2~\cite{PDL_D29_2026}]
\label{thm:val_unique}
Under axioms C1--C4, $\rval = 930$ is the unique solution for the valence content of a
proton-type closure: it is the only value satisfying charge $+1$, divisibility by
$R_e = 6$, and the $(A)\wedge(B)$ structural requirement. Consequently, the valence
sector admits exactly \emph{one} $(A)\wedge(B)$-stable configuration.
\end{theorem}

\subsection{Sea exclusion}

\begin{proposition}[Sea exclusion, D29/D31~\cite{PDL_D29_2026,PDL_D31_2026}]
\label{prop:sea_excl}
The sea sector $\Gsea$ satisfies $(A)\wedge(B)$ with probability $1/4$ per
independent pulsation cycle. Over $N$ stationary cycles the probability of maintaining
$(A)\wedge(B)$ is $(1/4)^N \to 0$. The time-averaged stable coupling amplitude
from $\Gsea$ therefore vanishes in the stationary regime.
\end{proposition}

\subsection{Gate~3 and the active surface}

\begin{theorem}[Gate~3, D31/D36~\cite{PDL_D31_2026,PDL_D36_2026}]
\label{thm:gate3}
The effective gravitational coupling of an ensemble of $N$ nucleons satisfies
\begin{equation}
  \Geff(N) = \sig(N)\cdot\GPDL,
  \qquad
  \sig(N) = 1-(1-\kap)^N,
  \quad \kap = \frac{\Rsurf}{\Rtot} = \frac{310\vp}{11017} \in \mathbb{Q}(\sqrt{5}),
  \label{eq:gate3}
\end{equation}
proved under hypothesis H1 ($\kap = \Rsurf/\Rtot$ as engagement fraction)
and the independence of gravitational engagements (proved algebraically in D36).
\end{theorem}

% ══════════════════════════════════════════════════════════════════════════════
\section{BH-1: surface locality of relational entropy}
\label{sec:bh1}

\subsection{Relational entropy of a PDL closure}

We define the relational entropy of a \PDL\ closure $\mathcal{C}$ from the perspective
of an external observer, i.e., an entity that can probe $\mathcal{C}$ only through
the mixed triangles it forms with an external $\Kfour$ block.

\begin{definition}[Externally accessible microstates]
\label{def:microstates}
A \emph{microstate} of the \PDL\ closure $\mathcal{C}$ is an assignment of signs
$s \in \{+1,-1\}$ to each of its $\Rtot$ relations consistent with the \PDL\ axioms
(coherence C2, closure C3, pulsation C4). A microstate is
\emph{externally accessible} if it can be distinguished by an external $\Kfour$ block
through the stable coupling criterion $(A)\wedge(B)$. The number of externally
accessible microstates is denoted $\Omsurf(\mathcal{C})$.
\end{definition}

\begin{definition}[Relational entropy]
\label{def:entropy}
The \emph{relational entropy} of $\mathcal{C}$ is
\begin{equation}
  \SPDL(\mathcal{C}) = k_B \ln\,\Omsurf(\mathcal{C}).
  \label{eq:Spdl}
\end{equation}
\end{definition}

\subsection{The main proposition}

\begin{proposition}[BH-1: surface locality of relational entropy]
\label{prop:bh1}
Let $\mathcal{C}$ be a \PDL\ proton closure satisfying axioms C1--C4. Then:
\begin{enumerate}[label=(\roman*)]
  \item The valence sector $\Gval$ contributes exactly one externally accessible
        microstate: $\Omega_{\mathrm{val}} = 1$, hence $S_{\mathrm{val}} = 0$.
  \item The relational entropy of $\mathcal{C}$ is entirely carried by the active
        surface: $\SPDL(\mathcal{C}) = k_B\ln\,\Omega_{\mathrm{surf}}$ where
        $\Omega_{\mathrm{surf}}$ counts configurations of the surface sector alone.
  \item The relational entropy satisfies the area law:
        \begin{equation}
          \SPDL \;\propto\; \Rsurf, \label{eq:area_law}
        \end{equation}
        not $\SPDL \propto \Rtot$ (volume law).
\end{enumerate}
\end{proposition}

\begin{proof}
\textbf{Part (i).}
By Theorem~\ref{thm:val_unique} (D29 Lemma~2), the valence sector $\Gval$ occupies a
\emph{unique} configuration satisfying $(A)\wedge(B)$: there is exactly one way to
arrange $\rval = 930$ stationary relations consistently with charge $+1$,
divisibility by $R_e$, and the stability criterion. For any external $\Kfour$ block,
the stable coupling depends on the cross-product of $\Kfour$ signs with the proton
signs via condition~(A). Since the valence configuration is unique, its coupling
signature to $\Kfour$ is fixed: no second valence configuration exists that would
produce a distinguishable response. Hence $\Omega_{\mathrm{val}} = 1$ and
$S_{\mathrm{val}} = k_B\ln 1 = 0$.

\textbf{Part (ii).}
By Proposition~\ref{prop:sea_excl} (D29/D31), $\Gsea$ is excluded from the
stationary $(A)\wedge(B)$-stable coupling in the limit $N\to\infty$. The stationary
coupling amplitude is therefore determined by $\Gval$ alone in the valence sector
and by the active surface $\Rsurf$ in the interface sector. Since the valence
contributes $\Omega_{\mathrm{val}} = 1$ (Part~(i)), the total externally accessible
count factorises as
\begin{equation}
  \Omsurf(\mathcal{C}) = \Omega_{\mathrm{val}}\cdot\Omega_{\mathrm{surf}} = \Omega_{\mathrm{surf}},
\end{equation}
and the relational entropy reduces to $\SPDL = k_B\ln\,\Omega_{\mathrm{surf}}$.

\textbf{Part (iii).}
The active surface $\Rsurf = 310\vp$ is a finite set of relations. Its configuration
space has at most $2^{\Rsurf}$ elements, and at least one accessible configuration
(the proton surface in its ground state). For a uniform prior over accessible
surface configurations, $\Omega_{\mathrm{surf}} = f(\Rsurf)$ for some strictly
monotone function of $\Rsurf$ alone, giving $\SPDL \propto \Rsurf$. Volume-law
scaling $\SPDL \propto \Rtot$ would require $\Omega_{\mathrm{val}} > 1$, which
contradicts Part~(i).
\end{proof}

\begin{remark}[Physical interpretation]
\label{rem:interpretation}
Proposition~BH-1 states that a \PDL\ closure is a \emph{surface-information object}:
its full interior --- the $\rval = 930$ valence relations carrying the proton's mass
and charge structure --- is, from outside, in a definite unique state and contributes
zero entropy. All thermodynamic uncertainty resides at the boundary. This is not a
consequence of a holographic postulate added to \PDL; it is a theorem of the four
\PDL\ axioms together with the uniqueness result of D29.

This mirrors precisely the Bekenstein--Hawking picture for black holes, where the
interior geometry is inaccessible and the entropy is localised on the event horizon.
The \PDL\ analogy is structural, not metaphorical: the unique-state interior of the
\PDL\ proton plays the role of the black hole interior; the active surface $\Rsurf$
plays the role of the horizon.
\end{remark}

\begin{remark}[Direction of the closure]
\label{rem:direction}
As noted in the main \PDL\ paper~\cite{PDL_Main_2026}, the proton and the black hole
are ``closure objects in opposite directions'': the black hole draws coherence inward
toward its horizon, whilst the proton projects coherence outward through its active
surface. Proposition~BH-1 establishes that \emph{both} objects share the same
entropy localisation law: $S \propto \text{surface}$, regardless of direction.
\end{remark}

% ══════════════════════════════════════════════════════════════════════════════
\section{Connection to Bekenstein--Hawking at macroscopic scale}
\label{sec:macro}

\subsection{From relational surface to geometric area}

At macroscopic scale, the active surface $\Rsurf$ per nucleon maps to the
gravitational engagement fraction $\kap = \Rsurf/\Rtot$. For an ensemble of $N$
nucleons the effective engaged surface is $\sig(N) = 1-(1-\kap)^N$ (Gate~3,
Theorem~\ref{thm:gate3}), and the effective gravitational coupling is
$\Geff(N) = \sig(N)\cdot\GPDL$. The Schwarzschild radius of a black hole of
mass $M$ in the \PDL\ framework is therefore
\begin{equation}
  r_s(N,M) = \frac{2\,\Geff(N)\,M}{c^2} = \frac{2\,\sig(N)\,\GPDL\,M}{c^2},
  \label{eq:rs_N}
\end{equation}
and the horizon area scales as $A = 4\pi r_s^2 \propto \sig(N)^2\,M^2$.

\subsection{PDL Bekenstein--Hawking formula}

Substituting $G \to \Geff(N)$ in equation~\eqref{eq:BH}:
\begin{equation}
  \SBH^{\PDL}(N,M) = \frac{k_B c^3 A}{4\,\Geff(N)\,\hb}
  = \frac{\pi k_B c^3\,r_s(N,M)^2}{\Geff(N)\,\hb}
  = \frac{4\pi k_B\,\Geff(N)\,M^2}{\hb\,c}.
  \label{eq:SBH_PDL}
\end{equation}
For the local universe ($N = N_{\mathrm{loc}} \approx 120$, $\sig(120) = 0.9963$),
equation~\eqref{eq:SBH_PDL} recovers standard Bekenstein--Hawking to within $0.37\%$.
For the CMB epoch ($N_{\mathrm{CMB}} = 40$, $\sig(40) = 0.8449$), the entropy of a
black hole of fixed mass $M$ was smaller by a factor $\sig(40)/\sig(120) = 0.848$:
black holes in the primordial universe had $\sim 15\%$ less entropy per unit mass
than their local counterparts.

\subsection{The area law as a theorem}

\begin{corollary}[PDL area law]
\label{cor:area_law}
Under Proposition~BH-1 and the identification of $\Rsurf$ with the physical horizon
area via $\Geff(N)$ (Gate~3, Theorem~\ref{thm:gate3}), the Bekenstein--Hawking area
law $\SBH \propto A$ is a structural consequence of the four \PDL\ axioms. It is not
an independent postulate.
\end{corollary}

\begin{proof}
Proposition~BH-1 gives $\SPDL \propto \Rsurf$ (area law at the relational level).
Gate~3 identifies $\Rsurf$ with the gravitational engagement surface, which maps to
the geometric horizon area $A$ in the macroscopic limit via $\Geff(N) = \sig(N)\cdot\GPDL$.
Composing these two identifications yields $\SBH \propto A$.
\end{proof}

% ══════════════════════════════════════════════════════════════════════════════
\section{Open problems}
\label{sec:open}

\begin{openprob}[The coefficient $1/4$]
\label{op:quarter}
Corollary~\ref{cor:area_law} establishes the area law but not the numerical
coefficient. The Bekenstein--Hawking formula~\eqref{eq:BH} has coefficient $1/4$
in Planck units: $\SBH = k_B\,A/(4\lP^2)$.
In \PDL, the analogous coefficient from Proposition~BH-1 would be
\begin{equation}
  \SPDL = k_B\ln\,\Omega_{\mathrm{surf}} = k_B\,\frac{\Rsurf}{\lambda_{\mathrm{PDL}}},
  \label{eq:coeff}
\end{equation}
where $\lambda_{\mathrm{PDL}}$ is the \PDL\ unit of relational area per microstate
(analogous to $4\lP^2$). The quantitative prediction $\lambda_{\mathrm{PDL}} = 4\lP^2$
requires: (a) the algebraic proof of $\kap = \Rsurf/\Rtot$ from axioms C1--C4 (OP1
of D36), and (b) a counting argument showing that each Planck area $\lP^2$ corresponds
to exactly four relational units of $\Rsurf$. This is the principal open problem of
D37 and constitutes the next step towards a fully parameter-free \PDL\ derivation
of Bekenstein--Hawking.
\end{openprob}

\begin{openprob}[Counting of surface microstates]
\label{op:counting}
Proposition~BH-1(iii) establishes $\SPDL \propto \Rsurf$ under the assumption that
$\Omega_{\mathrm{surf}} = f(\Rsurf)$ for a strictly monotone function. A precise
count of the surface configurations consistent with the \PDL\ axioms (closure C3,
coherence C2) would make equation~\eqref{eq:coeff} explicit and determine the
proportionality constant analytically.
\end{openprob}

\begin{openprob}[Hawking radiation from \PDL\ pulsation]
\label{op:hawking}
Equation~\eqref{eq:SBH_PDL} assigns an entropy to a \PDL\ black hole but does not
derive a temperature or an evaporation mechanism. Hawking radiation corresponds,
in the \PDL\ picture, to irreversible leakage of coherence through the active
surface: $\dot{S} = -k_B\,\dot{\Omega}_{\mathrm{surf}}/\Omega_{\mathrm{surf}}$.
Deriving the Hawking temperature $T_H = \hb c^3/(8\pi G_{\mathrm{eff}}(N)\,M\,k_B)$
from the pulsation dynamics of the \PDL\ closure is an open problem at this stage.
\end{openprob}

% ══════════════════════════════════════════════════════════════════════════════
\section{Epistemic status}
\label{sec:epistemic}

\begin{table}[ht]
\centering
\caption{Epistemic status of the results of D37.}
\label{tab:epistemic}
\medskip
\begin{tabular}{>{\raggedright}p{6.2cm}p{3.5cm}p{3.8cm}}
\toprule
Statement & Status & Dependency \\
\midrule
$\Omega_{\mathrm{val}} = 1$ (valence has one accessible state)
  & Theorem
  & D29 Lemma~2 (proved) \\[4pt]
$\SPDL \propto \Rsurf$ (area law at relational level)
  & Theorem (BH-1)
  & D29, D31 \\[4pt]
$\Geff(N) = \sig(N)\cdot\GPDL$
  & Theorem under H1
  & D31/D36, H1 = OP1 D36 \\[4pt]
$\SBH^{\PDL} \propto A$ (macroscopic area law)
  & Theorem under H1
  & BH-1 + Gate~3 \\[4pt]
Coefficient $1/4$ is a \PDL\ theorem
  & Open problem (OP1 D37)
  & OP1 D36 + microstate count \\[4pt]
Hawking temperature from \PDL\ pulsation
  & Open problem (OP3 D37)
  & Pulsation dynamics \\
\bottomrule
\end{tabular}
\end{table}

% ══════════════════════════════════════════════════════════════════════════════
\section{Conclusion}

This paper has established one main result and one corollary within the \PDL\
programme.

\emph{Main result (Proposition~BH-1).} The information-theoretic entropy of a \PDL\
closure is entirely carried by its active surface $\Rsurf$, not by its total
relational budget~$\Rtot$. This follows from the unique-state property of the valence
sector (D29 Lemma~2): since the interior is frozen in a single $(A)\wedge(B)$-stable
configuration, it contributes zero entropy to any external observer. The area law
$S \propto \Rsurf$ is therefore an unconditional theorem of the \PDL\ axioms.

\emph{Corollary (PDL area law).} Under Gate~3 ($\Geff(N) = \sig(N)\cdot\GPDL$),
the macroscopic Bekenstein--Hawking area law $\SBH \propto A$ is a structural
consequence of the \PDL\ axioms, not an independent postulate. The coefficient
$1/4$ remains an open quantitative problem (OP1 of D37), conditional on the
algebraic proof of $\kap = \Rsurf/\Rtot$ (OP1 of D36).

The structural identity between the entropy-area law of black holes and the
surface-locality of relational entropy in the \PDL\ proton is not metaphorical:
both follow from the same organisational principle --- information on the boundary,
inaccessible unique interior --- applied to two radically different scales of closure.
Whether the \PDL\ framework can extend this identity to derive Hawking radiation and
the cosmological constant from the same combinatorial object $\kap$ remains the
principal open frontier of the programme.

% ══════════════════════════════════════════════════════════════════════════════
\bibliography{Pdl_refs_d37.bib}

\end{document}